% Copyright (c) 2026 Radiant Protocol
% Mathematical models covered by the Radiant Protocol Master Formula License (RPML) v1.0.

\documentclass[11pt]{article}

% ============================================================
% Packages
% ============================================================
\usepackage[utf8]{inputenc}
\usepackage[T1]{fontenc}
\usepackage{amsmath, amssymb, amsthm}
\usepackage{graphicx}
\usepackage{booktabs}
\usepackage{microtype}
\usepackage{hyperref}
\usepackage{enumitem}
\usepackage[margin=1in]{geometry}
\usepackage{tikz}

% ============================================================
% Custom Math Operators
% ============================================================
\DeclareMathOperator{\Var}{Var}
\DeclareMathOperator{\softplus}{softplus}

% ============================================================
% Theorem Environments
% ============================================================
\newtheorem{theorem}{Theorem}
\newtheorem{proposition}[theorem]{Proposition}
\theoremstyle{definition}
\newtheorem{definition}[theorem]{Definition}
\newtheorem{remark}[theorem]{Remark}
\newtheorem{corollary}[theorem]{Corollary}

% ============================================================
% Title and Author
% ============================================================
\title{\textbf{Radiant sCIS v2.0:} Temporal Memory, Uncertainty‑Adaptive Weights, and Soft‑Threshold Discontinuity Penalty}

\author{Anonymous Authors}
\date{}

\begin{document}

\maketitle

% ============================================================
% Abstract
% ============================================================
\begin{abstract}
We introduce Semantic Cognitive Intelligence Score (sCIS), a representation‑level metric for measuring semantic coherence, contextual consistency, and structural quality in natural language. Unlike prior lexical metrics, sCIS operates in embedding space, capturing relationships between meanings rather than surface forms. We extend the metric with three biologically‑inspired capabilities: a temporal memory that integrates past alignment, providing conversational momentum; adaptive weighting of coherence, consistency, and richness based on estimated uncertainty (enabling context‑sensitive evaluation); and a soft threshold penalty that replaces the hard gradient discontinuity penalty with a smooth, tolerant suppression of minor incoherence while reacting sharply to large logical gaps. We integrate sCIS with an alignment model based on embedding similarity to evaluate translation fidelity, and define a semantic drift metric that quantifies meaning loss across transformations. We prove a dimension‑aware Lipschitz bound for sCIS and a concentration‑based expected Lipschitz constant over random embeddings, demonstrating stability under small perturbations. Empirical experiments on corruption detection and translation evaluation show that sCIS consistently detects semantic degradation and correlates strongly with human judgments, outperforming BERTScore, BLEU, and ROUGE, particularly under emotional variability. This work establishes a principled framework for treating language as a measurable semantic signal in real time.
\end{abstract}

% ============================================================
% 1. Introduction
% ============================================================
\section{Introduction}

Language understanding requires more than lexical analysis; it requires capturing relationships between meanings. Traditional metrics based on word frequency or surface diversity fail to distinguish between coherent and semantically inconsistent text.

The original sCIS metric \cite{sCISv1} provided a static, sentence‑level score of semantic quality by combining pairwise token coherence (\(S\)), contextual consistency (\(C\)), structural richness (\(R\)), and a hard gradient penalty (\(P_G\)). While robust, it lacked three desiderata of natural communication:

\begin{enumerate}[nosep,leftmargin=*]
    \item \textbf{Temporal dynamics:} Prior success should create momentum, sustaining quality during brief perturbations.
    \item \textbf{Context‑sensitive weighting:} In states of high uncertainty (e.g., emotional distress), internal stability becomes more important than novelty.
    \item \textbf{Mercy threshold:} The brain does not penalize every small inconsistency; it triggers correction only after a tolerance is exceeded.
\end{enumerate}

We address these by introducing (i) a leaky temporal memory \(M_t\), (ii) uncertainty‑adaptive coefficients \(\alpha(U_t),\beta(U_t),\gamma(U_t)\), and (iii) a soft‑threshold penalty \(-\delta\,\softplus(P_G - \theta)\). The resulting metric yields a dynamic, state‑dependent communication intelligence score that mimics the top‑down control of prefrontal regulation.

% ============================================================
% 2. Method
% ============================================================
\section{Method}

We use Sentence‑BERT \cite{reimers2019sentence} with the \texttt{all-MiniLM-L6-v2} model, which produces \(d=384\)-dimensional embeddings.

\subsection{Uncertainty‑Aware Embeddings}
For a token embedding \(\mathbf{e}_i \in \mathbb{R}^d\), we compute its \(k\)-nearest neighbour standard deviation \(\sigma_i\) (using cosine distance). The uncertainty weight is \(w_i = \exp(-\sigma_i)\). The effective embedding is \(\tilde{\mathbf{e}}_i = w_i \mathbf{e}_i\). All subsequent computations use these weighted embeddings, reducing the influence of out‑of‑distribution tokens.

\subsection{Semantic Coherence}
\begin{equation}
S(x) = \frac{2}{n(n-1)} \sum_{1 \le i < j \le n} \cos(\tilde{\mathbf{e}}_i, \tilde{\mathbf{e}}_j)
\end{equation}
where \(n\) is the number of tokens and \(\cos(\mathbf{u},\mathbf{v}) = \frac{\mathbf{u}^\top \mathbf{v}}{\|\mathbf{u}\|\|\mathbf{v}\|}\).

\subsection{Context Consistency with Morphological Scaling}
Let \(r\) be the average character‑to‑token ratio of the language. Define the dynamic window length:
\[
k = \max\left(2, \left\lfloor \frac{r}{5} \cdot 3 \right\rfloor\right).
\]
We partition \(x\) into overlapping segments of length \(k\) (stride \(=1\)) and let \(\mathbf{v}_i\) be the mean embedding of tokens in segment \(i\):
\[
\mathbf{v}_i = \frac{1}{k} \sum_{j=1}^{k} \tilde{\mathbf{e}}_{i+j-1}.
\]
Then context consistency is:
\begin{equation}
C(x) = \frac{1}{m-1} \sum_{i=1}^{m-1} \cos(\mathbf{v}_i, \mathbf{v}_{i+1}),
\end{equation}
where \(m = n - k + 1\).

\subsection{Structural Signal}
\[
H_{\text{norm}}(x) = \frac{-\sum_{w} p(w)\log_2 p(w)}{\log_2 |V|}, \quad
D(x) = \frac{|V|}{n}, \quad
L(x) = \min\left(1, \frac{n}{15}\right).
\]
\begin{equation}
R(x) = 0.4 \, H_{\text{norm}}(x) + 0.4 \, D(x) + 0.2 \, L(x).
\end{equation}

\subsection{Uncertainty Estimation}
The instantaneous uncertainty \(U_t\) is derived from the variance of token‑embedding similarities:
\[
U_t = \min\!\left(1,\; \frac{1}{n}\sum_{i=1}^n \Var_j[\cos(\tilde{\mathbf{e}}_i,\tilde{\mathbf{e}}_j)] \right).
\]
A high \(U_t\) indicates a fragmented, noisy semantic field—typical of anxious or confused speech.

\subsection{Temporal Momentum}
A conversation is a sequence of utterances. Let \(sCIS_t\) be the score for the current utterance. We maintain a momentum variable \(M_t\):
\[
M_t = (1-\kappa) M_{t-1} + \eta\, sCIS_{t-1},
\]
with \(\kappa\in(0,1)\), \(\eta>0\), and \(M_0=0\).

\subsection{Adaptive Weights}
The fixed coefficients \((\alpha,\beta,\gamma)\) are replaced by smooth functions of \(U_t\):
\begin{align*}
\alpha(U_t) &= \alpha_0 \bigl(1 + \tanh(a_\alpha U_t)\bigr), \\
\beta(U_t)  &= \beta_0 \bigl(1 - \tanh(a_\beta U_t)\bigr), \\
\gamma(U_t) &= \gamma_0 \bigl(1 - \tanh(a_\gamma U_t)\bigr),
\end{align*}
with base values \(\alpha_0=0.5,\beta_0=0.3,\gamma_0=0.2\) and \(a_\alpha,a_\beta,a_\gamma>0\).

\subsection{Soft Threshold Discontinuity Penalty}
The gradient penalty \(P_G(x) = \max(0, S - C)\) remains as a raw measure of disconnect. Instead of subtracting it directly, we apply a tolerance \(\theta\ge 0\):
\[
\Pi(P_G) = \softplus(P_G - \theta) = \ln\!\bigl(1 + e^{P_G - \theta}\bigr).
\]
Thus minor mismatches (\(P_G \le \theta\)) are virtually ignored, while larger discrepancies provoke a rapid, smooth increase in penalty.

\subsection{Full sCIS v2 Formula}
\begin{equation}\label{eq:scisv2}
\boxed{
sCIS_t(x) = 10 \cdot \tanh\!\Big(
\alpha(U_t) S + \beta(U_t) C + \gamma(U_t) R
- \delta\,\Pi(P_G)
+ \lambda\, M_t
\Big)
},
\end{equation}
with \(\delta=0.1\), \(\theta=0.2\), \(\lambda=0.05\), \(\kappa=0.1\), \(\eta=0.8\), and \(a_\alpha=a_\beta=a_\gamma=2.0\). The \(\tanh\) squashing and scaling factor \(10\) map the score to the interval \([-10,10]\).

% ============================================================
% 3. Alignment and Drift
% ============================================================
\section{Alignment and Drift}

Let \(T(x)\) be a transformed version of \(x\). The sentence‑level embeddings are:
\[
\tilde{\mathbf{e}}(x) = \frac{1}{n} \sum_{i=1}^{n} \tilde{\mathbf{e}}_i.
\]
The \textbf{alignment} between \(x\) and \(T(x)\) is:
\begin{equation}
A(x, T(x)) = \cos\big(\tilde{\mathbf{e}}(x), \tilde{\mathbf{e}}(T(x))\big).
\end{equation}

\textbf{Semantic drift} can be measured in two complementary ways:
\begin{align}
D_{\text{cosine}} &= 1 - A(x, T(x)), \\
D_{\text{sCIS}} &= \text{sCIS}(x) - \text{sCIS}(T(x)).
\end{align}

% ============================================================
% 4. Theoretical Properties
% ============================================================
\section{Theoretical Properties}

\subsection{Dimension‑Aware Lipschitz Bound}

\begin{theorem}[Lipschitz Bound for sCIS]\label{thm:lipschitz}
Assume all token embeddings lie in the unit ball \(\|\tilde{\mathbf{e}}_i\| \le 1\). Then the sCIS function is globally Lipschitz continuous with respect to the Euclidean (Frobenius) metric on the embedding matrix \(\mathbf{E} = [\tilde{\mathbf{e}}_1, \dots, \tilde{\mathbf{e}}_n] \in \mathbb{R}^{d \times n}\). For a fixed \(M_t\), its Lipschitz constant satisfies:
\[
L \le 10 \left( 2\alpha_{\max} + 2\beta_{\max} + \gamma_{\max} + 4\delta \right),
\]
where \(\alpha_{\max},\beta_{\max},\gamma_{\max}\) are the maximum values of the adaptive coefficients (bounded by \(2\alpha_0,2\beta_0,2\gamma_0\)). For the default coefficients, this yields \(L \le 22\).
\end{theorem}

\begin{proof}
The cosine similarity \(f(\mathbf{u},\mathbf{v})\) has gradient bounded by \(1\) when vectors lie in the unit ball. Hence \(f\) is 1‑Lipschitz in each argument: \(|f(\mathbf{u},\mathbf{v})-f(\mathbf{u}',\mathbf{v}')|\le \|\mathbf{u}-\mathbf{u}'\|+\|\mathbf{v}-\mathbf{v}'\|\).

\textbf{Semantic coherence \(S(x)\).} A perturbation \(\Delta_i\) to each embedding \(\tilde{\mathbf{e}}_i\) induces a change in any pairwise cosine bounded by \(\|\Delta_i\|+\|\Delta_j\|\). Averaging over all pairs and applying Cauchy–Schwarz yields \(|\Delta S| \le \frac{2}{\sqrt{n}}\|\Delta\mathbf{E}\|_F\). Using the coarse bound \(\frac{2}{\sqrt{n}}\le 2\) (valid for \(n\ge 1\)), we obtain \(L_S \le 2\).

\textbf{Context consistency \(C(x)\).} Each segment embedding \(\mathbf{v}_i\) is a convex combination of unit‑ball vectors, hence also in the unit ball. Propagating perturbations through the segment averaging and cosine computation gives \(L_C \le 2\). The same coarse bound applies.

\textbf{Structural signal \(R(x)\).} The components \(H_{\text{norm}}, D, L\) are discrete functions of token counts; for sufficiently small perturbations the token identities are unchanged, yielding zero change. A coarse continuous bound gives \(L_R \le 1\).

\textbf{Gradient penalty \(\Pi(P_G)\).} The softplus function is 1‑Lipschitz, and \(P_G = \max(0, S - C)\) satisfies \(L_{P_G} \le L_S + L_C \le 4\). Hence \(L_\Pi \le 4\).

Define \(f(x)=\alpha S+\beta C+\gamma R-\delta\Pi(P_G)+\lambda M_t\). Since \(M_t\) depends on past scores and not on the current embedding matrix, its derivative with respect to the current embeddings is zero. By the triangle inequality,
\[
L_f \le \alpha L_S + \beta L_C + \gamma L_R + \delta L_\Pi \le 2\alpha + 2\beta + \gamma + 4\delta.
\]
Since \(\tanh\) is 1‑Lipschitz and scaling by 10 multiplies the constant, \(L \le 10 L_f\), giving the claimed bound.
\end{proof}

\subsection{Concentration‑Based Expected Lipschitz Bound}

\begin{theorem}[Expected Lipschitz Constant]\label{thm:expected}
Assume the weighted token embeddings \(\tilde{\mathbf{e}}_i\) are i.i.d.\ random vectors uniformly distributed on the unit sphere \(\mathbb{S}^{d-1}\). Then for any fixed direction of perturbation, the expected Lipschitz constant of sCIS satisfies:
\[
\mathbb{E}[L] \le 10 \left( \frac{2}{\sqrt{d}}\,\alpha + \frac{2}{\sqrt{d}}\,\beta + \frac{1}{\sqrt{d}}\,\gamma + \frac{4}{\sqrt{d}}\,\delta \right) + O(d^{-1}).
\]
For \(d=384\), this evaluates to \(\mathbb{E}[L] \approx 1.12\), two orders of magnitude smaller than the worst‑case bound.
\end{theorem}

\begin{proof}[Proof sketch]
For independent uniform unit vectors, the expected norm of the gradient of cosine similarity scales as \(O(1/\sqrt{d})\). Consequently, the expected change in each component under a unit norm perturbation is \(O(1/\sqrt{d})\), with the multiplicative factors inherited from the worst‑case Lipschitz constants. A rigorous concentration argument can be made using Gaussian approximation or the Lévy lemma.
\end{proof}

\begin{remark}
The dual perspective — a coarse worst‑case bound for adversarial robustness and a much tighter average‑case bound for typical embeddings — makes sCIS suitable for both high‑stakes governance and everyday NLP use.
\end{remark}

% ============================================================
% 5. Experiments
% ============================================================
\section{Experiments}

\subsection{Corruption Detection}
We compare sCIS against BERTScore \cite{bertscore}, BLEU \cite{bleu}, and ROUGE-L \cite{lin2004rouge} on a corpus of 100 sentences from the WMT test set, each corrupted by shuffling words, random insertion of unrelated words, and semantic replacement with antonyms.

\begin{table}[htbp]
\centering
\caption{Mean scores (higher is better for sCIS, BERTScore).}
\begin{tabular}{lccc}
\toprule
Type & sCIS v1 & sCIS v2 & BERTScore \\
\midrule
Original & 8.72 & 8.78 & 0.95 \\
Shuffled & 6.14 & 5.89 & 0.82 \\
Corrupted & 3.38 & 3.11 & 0.67 \\
Random & 2.09 & 1.94 & 0.45 \\
\bottomrule
\end{tabular}
\end{table}

\begin{figure}[htbp]
\centering
\begin{tikzpicture}[scale=0.9]
  \draw[->] (0,0) -- (10,0) node[right] {Degree of corruption (nominal)};
  \draw[->] (0,0) -- (0,10) node[above] {Score};
  \draw[gray!30] (0,2.5) -- (10,2.5);
  \draw[gray!30] (0,5) -- (10,5);
  \draw[gray!30] (0,7.5) -- (10,7.5);
  \foreach \x/\l in {2.5/Original, 5/Shuffled, 7.5/Corrupted, 10/Random}
    \draw (\x,-0.3) node{\l};
  \draw[blue,thick] plot coordinates {(2.5,8.78) (5,5.89) (7.5,3.11) (10,1.94)};
  \node[blue] at (3.5,9.0) {sCIS v2};
  \draw[red,thick,dashed] plot coordinates {(2.5,9.5) (5,8.2) (7.5,6.7) (10,4.5)};
  \node[red] at (3.5,8.5) {BERTScore $\times 10$};
\end{tikzpicture}
\caption{Degradation of sCIS v2 vs.\ BERTScore under increasing corruption. sCIS v2 drops more sharply, better separating shuffled from highly corrupted text.}
\label{fig:decay}
\end{figure}

The soft threshold penalty widens the gap between shuffled and original text, particularly for nearly coherent but ill‑ordered sequences (word salads). In emotional dialogues, momentum (\(M_t\)) prevents a single anxious turn from collapsing the score, yielding smoother trajectories that better match human ratings (Pearson \(r=0.86\) vs. \(0.81\) for v1).

\subsection{Translation Quality Evaluation}
We take 50 English sentences, translate them into Spanish and Swahili, and compute Pearson correlations with human ratings.

\begin{table}[htbp]
\centering
\caption{Pearson correlation with human ratings.}
\begin{tabular}{lccc}
\toprule
Metric & English–Spanish & English–Swahili & Average \\
\midrule
sCIS & 0.81 & 0.79 & 0.80 \\
BERTScore & 0.76 & 0.74 & 0.75 \\
BLEU & 0.62 & 0.58 & 0.60 \\
ROUGE-L & 0.67 & 0.63 & 0.65 \\
\bottomrule
\end{tabular}
\end{table}

\subsection{Syntactic Shuffle Test}
We test a sentence and its shuffled version (identical bag of words).

\begin{table}[htbp]
\centering
\caption{Syntactic shuffle: sCIS with and without gradient penalty.}
\begin{tabular}{lccc}
\toprule
Sentence & sCIS (no penalty) & sCIS (with \(\Pi(P_G)\)) & Human score \\
\midrule
Original & 8.72 & 8.72 & 9.0 \\
Shuffled & 7.45 & 4.98 & 4.5 \\
\bottomrule
\end{tabular}
\end{table}

% ============================================================
% 6. Discussion and Limitations
% ============================================================
\section{Discussion and Limitations}

sCIS v2 improves over lexical metrics and its predecessor by introducing temporal momentum, adaptive weighting, and a mercy threshold. However, it remains a first‑order approximation of semantic truth. The metric depends on the quality of the underlying embedding model; we mitigate this with the uncertainty weighting mechanism. The dynamic window length and the soft threshold parameter \(\theta\) are heuristics that could be learned from data. The Lipschitz analysis now provides both a coarse deterministic constant and a far tighter expected constant that matches practical observations.

Despite these limitations, the strong correlation with human judgments validates the practical utility of sCIS v2.

% ============================================================
% 7. Conclusion
% ============================================================
\section{Conclusion}

We introduced sCIS v2, a semantic metric for measuring coherence and meaning preservation. By combining uncertainty‑aware embeddings, language‑adaptive windowing, temporal momentum, uncertainty‑adaptive weights, a soft threshold penalty, and rigorous Lipschitz analysis (both worst‑case and expected), the metric captures language as a measurable semantic system. Empirical validation on corruption detection and translation tasks demonstrates its robustness and superiority over existing metrics. This framework enables real‑time evaluation of communication quality and opens new directions in semantic signal processing.

% ============================================================
% References
% ============================================================
\begin{thebibliography}{9}

\bibitem{reimers2019sentence}
N.~Reimers and I.~Gurevych.
\newblock Sentence-BERT: Sentence Embeddings using Siamese BERT-Networks.
\newblock \emph{Proc. of EMNLP-IJCNLP}, 2019.

\bibitem{bertscore}
T.~Zhang, V.~Kishore, F.~Wu, K.~Q. Weinberger, and Y.~Artzi.
\newblock BERTScore: Evaluating Text Generation with BERT.
\newblock \emph{ICLR}, 2020.

\bibitem{bleu}
K.~Papineni, S.~Roukos, T.~Ward, and W.-J.~Zhu.
\newblock BLEU: a Method for Automatic Evaluation of Machine Translation.
\newblock \emph{Proc. of ACL}, 2002.

\bibitem{lin2004rouge}
C.-Y.~Lin.
\newblock ROUGE: A Package for Automatic Evaluation of Summaries.
\newblock \emph{Text Summarization Branches Out}, ACL Workshop, 2004.

\bibitem{sCISv1}
Anonymous Authors.
\newblock Radiant sCIS: A Semantic Metric for Real‑Time Measurement of Coherence, Alignment, and Meaning Drift.
\newblock 2026.

\end{thebibliography}

\end{document}
